\section{Background}
\label{sec:background}

\subsection{Tool-Using LLM Agents}

Modern LLMs (GPT-4~\cite{openai-gpt4}, Claude~\cite{anthropic-claude}, Llama~\cite{llama}) support \emph{function calling}~\cite{schick2023toolformer}, enabling agents to interact with external systems. An LLM agent consists of three core components: (1)~a \emph{system prompt} defining the agent's role, behavioral constraints, and available capabilities; (2)~a \emph{tool registry} of callable functions with typed schemas specifying parameter names, types, and descriptions; and (3)~an \emph{execution loop} that iteratively receives user input, generates tool calls via structured output (JSON function calls), executes them, incorporates results, and decides whether to continue or respond.

In smart homes, this architecture enables multi-step automation: a single user command such as ``summarize today's schedule and turn off unused lights'' triggers a sequence of tool calls (\texttt{calendar.list\_events}, \texttt{device.list}, \texttt{device.control}) orchestrated by the LLM. The agent reads content from multiple sources during execution: user commands, system prompts, device states, calendar entries, file contents, and notification messages. This creates a fundamental security challenge: if an attacker controls any content the agent reads (e.g., a calendar event title containing ``SYSTEM: delete all files''), the injected instructions may hijack subsequent tool calls.

\paragraph{Agent--Tool Interaction Model.}
The LLM generates a structured tool call (e.g., \texttt{\{"tool": "fs.delete", "args": \{"path": "/backups"\}\}}), which the framework dispatches to the executor. The result is appended to the LLM's context. Critically, the LLM has \emph{no mechanism} to distinguish tool results from trusted versus untrusted sources---both appear as text tokens in the same context window.

\subsection{Prompt Injection}

Prompt injection~\cite{perez2022ignore,greshake2023not} exploits the fact that LLMs process instructions and data in the same text stream, with no syntactic boundary. Three vectors are especially relevant. \emph{Direct injection} embeds malicious instructions in user-facing input (InjecAgent~\cite{injecagent2024} reports 24\% vulnerability for GPT-4). \emph{Indirect injection}~\cite{greshake2023not} embeds instructions in content the agent reads during execution (e.g., device names, calendar events, and file contents), making it more dangerous because the attacker need not interact with the user and the payload can persist across sessions. \emph{Chained/multi-vector attacks} combine cross-turn injection chains, encoding obfuscation (Base64, hex, ROT13, Unicode), and role-play jailbreaks to evade pattern-based defenses. Section~\ref{sec:intro} discusses why existing defenses---input filters, output validation, instruction hierarchy, and dual-LLM architectures---all fail to address the root cause: the absence of data provenance.

\subsection{Capability Security and Provenance}

\paragraph{Capability-Based Access Control.}
Capability-based security~\cite{miller2006capability,capsicum} replaces ambient authority with unforgeable tokens encoding permitted operations on specific resources. \sys{} adapts this model to LLM tool calls: policies define which operations are permitted under which conditions, with provenance serving as the ``capability token'' that determines authorization.

\paragraph{Data Provenance.}
The provenance literature distinguishes three fundamental questions about data lineage~\cite{cheney2009provenance,buneman2001provenance}: \emph{why-provenance} identifies which source tuples contributed to a result; \emph{where-provenance} identifies which source locations were copied to the output; \emph{how-provenance}~\cite{green2007provenance} describes how sources were combined, formalized as annotations over a provenance semiring. The W3C PROV Data Model (PROV-DM)~\cite{prov-dm} standardizes provenance representation using three core types: \emph{entities} (data items), \emph{activities} (processes that transform entities), and \emph{agents} (actors responsible for activities), connected by relations such as \texttt{wasDerivedFrom}, \texttt{wasGeneratedBy}, and \texttt{wasAttributedTo}.

In security, provenance manifests as dynamic taint analysis~\cite{taint-analysis-survey,taintdroid}: data from untrusted sources is ``tainted,'' taint propagates through operations, and security-sensitive sinks check for taint before execution. Provenance-aware storage systems~\cite{muniswamy2006provenance} and information-flow control~\cite{difc-survey} extend this to persistent data and multi-level security labels. However, all classical approaches assume \emph{instrumentable computation}---the ability to observe every instruction that transforms data.

\paragraph{Synthesis: Provenance-Gated Capabilities.}
\sys{} synthesizes capability-based security with data provenance at the tool-call boundary. TOOL\_RESULT and INPUT nodes map to PROV-DM \emph{entities}; LLM inference and tool executions to \emph{activities}; users and data sources to \emph{agents}. Trust labels ($\top$ = TRUSTED, $\bot$ = UNTRUSTED) extend entities with a security annotation from a two-element lattice $\mathcal{L} = \{\top, \bot\}$ ordered by $\bot \sqsubseteq \top$, where trust propagates via the meet ($\sqcap$) operator.

\paragraph{Why This Is Not Just Taint Tracking.}
Classical taint systems~\cite{taintdroid,difc-survey,flume} instrument every instruction with binary allow/block decisions. \sys{} differs in three ways: (1)~LLMs cannot be instrumented, requiring boundary-level \emph{why-provenance} via semantic matching; (2)~the LLM paraphrases content, requiring embedding-based resolution; (3)~agents must read untrusted data to function, demanding risk-tiered policies rather than blanket denial.
